\subsection{Symplectic modular symmetry in heterotic string vacua: flavor, CP, and $R$-symmetries --- arXiv:2107.00487}

\textbf{Authors:} Keiya Ishiguro, Tatsuo Kobayashi, Hajime Otsuka

\paragraph{主な主張}
標準埋め込みをもつヘテロティックCalabi--Yauコンパクト化で、フレーバー・CP・$U(1)_R$ の変換を一般化シンプレクティック構造に統一し、orbifoldおよび滑らかなCalabi--Yauの具体例で非自明なフレーバー対称性を示した\cite{Ishiguro:2021ccl}。

\paragraph{新規性}
物質場とモジュライの対応を利用して、モジュライ空間のシンプレクティック変換から物質場のフレーバー変換とそのCP拡張を導いた。

\paragraph{位置付け}
モジュラーフレーバー対称性の幾何学的起源を、トーラスのorbifoldから一般のCalabi--Yauの例へ広げる研究である。

\paragraph{コメント（読書メモ）}
読書メモ：Calabi--Yauの幾何に自然に現れるシンプレクティック変換を、モジュラーフレーバー対称性と結びつける論文として読んだ。$E_6$ の $\mathbf{27}$・$\overline{\mathbf{27}}$ 物質場にもシンプレクティック変換を定義しているのか、という点は当時の確認事項として残す。物質場は $U(1)_R$ 電荷をもち、対称性の関係は $Sp(2h+2,\mathbb{C})\supset Sp(2h+2,\mathbb{Z})\times U(1)_R$、$GSp(2h+2,\mathbb{Z})\supset G_{\rm flavor}\rtimes\mathbb{Z}_2^{\rm CP}$ と整理した。CPと $U(1)_R$ を含む構造は $GSp(2h+2,\mathbb{C})$ の中で捉える。\par
標準埋め込みでは物質場とモジュライが対応するため、フレーバー構造はCalabi--Yau 3-foldのシンプレクティック構造に強く関係する、と理解した。CPはシンプレクティック群の外部自己同型として、$U(1)_R$ は正則3形式の回転として捉える（元メモの「3-form fluxの回転」は正則3形式の回転に修正）。orbifoldの例だけでなくCalabi--Yauの例でも、Yukawa結合のフレーバー対称性がCPを含む $G_{\rm flavor}\rtimes\mathbb{Z}_2^{\rm CP}$ に拡張される点に注目した。
